% Generated draft: requires mathematical and source review.
\documentclass[11pt]{article}
\usepackage[margin=25mm]{geometry}
\usepackage{fontspec}
\setmainfont{Latin Modern Roman}
\usepackage{amsmath,amssymb,mathtools,unicode-math}
\setmathfont{Latin Modern Math}
\usepackage{enumitem,needspace,xurl,hyperref}
\setlength{\parindent}{0pt}
\setlength{\parskip}{5pt}
\setlength{\emergencystretch}{3em}
\newcommand{\N}{\mathbb{N}}
\newcommand{\Z}{\mathbb{Z}}
\newcommand{\Q}{\mathbb{Q}}
\newcommand{\R}{\mathbb{R}}
\newcommand{\C}{\mathbb{C}}
\newcommand{\F}{\mathbb{F}}
\newcommand{\A}{\mathbb{A}}
\newcommand{\PP}{\mathbb{P}}
\newcommand{\E}{\mathbb{E}}
\newcommand{\eps}{\varepsilon}
\newcommand{\dd}{\,\mathrm{d}}
\begin{document}
\section*{UnsolvedMath ID: 800589}
\textbf{Source problem code:} catalogue entry
\subsection{Definitions and mathematical statement}
Let \(\N=\{1,2,\ldots\}\). The pseudointersection number \(\mathfrak p\) is
\[
\mathfrak p=\min\bigl\{|\mathcal A|:\mathcal A\subseteq[\N]^{\aleph_0},\ \text{every finite subfamily of }\mathcal A\text{ has infinite intersection, and no infinite }B\subseteq\N\text{ satisfies }B\setminus A\text{ finite for every }A\in\mathcal A\bigr\}.
\]
This is the conventional meaning of the otherwise undefined symbol \(\mathfrak p\) in the supplied record.

For a metric space \((Z,d)\), Hausdorff dimension zero means that for every \(s>0\), \(\delta>0\), and \(\varepsilon>0\), there is a countable cover \((A_k)_{k\geq1}\) of \(Z\) such that
\[
\operatorname{diam}_d(A_k)\leq\delta\quad(k\geq1),\qquad
\sum_{k\geq1}\bigl(\operatorname{diam}_d(A_k)\bigr)^s<\varepsilon.
\]
Write this condition as \(\dim_{\mathrm H}(Z,d)=0\).

\textbf{Question (source Problem 15).} Assume \(|X|<\mathfrak p\). Is it true that for every admissible space \(Y\),
\[
\dim_{\mathrm H}Y=0\quad\Longrightarrow\quad\dim_{\mathrm H}(X\times Y)=0?
\]
The supplied record does not specify the ambient class of \(X\) and \(Y\), their metrics, or the product metric. Consequently, ``admissible'' remains unresolved: the displayed question is a statement template, not an assertion that the quantifier ranges over all metric spaces. Hausdorff dimension requires these metric choices, so a fully specified formulation requires recovery of the intended setting.

\textbf{Supplied context.} Conventionally, an open \(\omega\)-cover of \(X\) is an open cover that does not contain \(X\) itself and contains a member covering each finite subset of \(X\). An open \(\gamma\)-cover is an infinite family of open sets such that each point of \(X\) belongs to all but finitely many members. Let \(\Omega\) and \(\Gamma\) denote these respective classes. The property \(\mathsf S_1(\Omega,\Gamma)\) means that from each sequence \((\mathcal U_n)_{n\geq1}\) of members of \(\Omega\), one can select \(U_n\in\mathcal U_n\) so that \(\{U_n:n\geq1\}\in\Gamma\).

The record defines an open \(n\)-cover \(\mathcal U\) by requiring that every \(F\subseteq X\) with \(|F|\leq n\) be contained in some member of \(\mathcal U\). It describes \(\mathsf S_1(\{\mathrm O_n\}_{n\in\N},\Gamma)\), called the strong \(\gamma\)-property and associated with reference labels [GM, strongdiags], as selection of \(U_n\in\mathcal U_n\) from each sequence of open \(n\)-covers to obtain a \(\gamma\)-cover. The supplied text actually writes \(\{\mathcal U_n:n\in\N\}\) as that resulting cover; the interpretation as \(\{U_n:n\in\N\}\) is a proposed correction requiring review. Whether the open \(n\)-covers must exclude \(X\) itself is also unstated.

For \(Y\) of Hausdorff dimension zero, the source asserts that \(\mathsf S_1(\Omega,\Gamma)\) for \(X\) does not imply Hausdorff dimension zero of \(X\times Y\), whereas the strong \(\gamma\)-property does imply it, citing [Hdim] for the latter assertion. These are supplied contextual claims in an unspecified metric setting, not independently verified results of this draft.
\subsection{Short English statement}
The question asks whether a set \(X\) of cardinality below the pseudointersection number \(\mathfrak p\) has a product of Hausdorff dimension zero with every eligible space \(Y\) of Hausdorff dimension zero. The intended spaces and metrics are missing. The registry labels the problem open; that label and the contextual claims are supplied information requiring mathematical review.
\subsection{Sources}
\begin{itemize}
\item Local UnsolvedMath registry: permanent ID \(800589\), exact source problem code \texttt{catalogue entry}, upstream ID \(800589\), revision \texttt{d371ee917ac9647fbb06cd6fc529deec6f3dda82}. The supplied record attributes the problem to Tsaban, with proposed year 2006, and labels its status \texttt{open}.
\item Supplied source list: Tsaban, \emph{Selection Principles and Special Sets of Reals: Open Problems} (2006), Problem 15. Supplied URL: \url{https://arxiv.org/abs/math/0609601}. The record reports access on 2026-10-16 and extraction by \texttt{arxiv\_latex}; these provenance details are supplied and were not independently verified.
\item Supplied attribution/context: [Hdim]. The record cites [Hdim] for the product assertion involving the strong \(\gamma\)-property and [GM, strongdiags] for that property's terminology or definition. No bibliographic expansions or separate links for these labels are supplied.
\end{itemize}
\end{document}
